\section{Discussion}
\label{sec:discussion}

\subsection{What the Structural Improvement Establishes}
The comparisons support a specific computational observation: an explicit reference-edge objective increases agreement with the selected reference edges under the archived implementation. Its effect persists across the available image pairs and several local reconstruction targets. This consistency makes the objective a useful, transparent control within the tested workflow. It does not, however, establish that the retained edges are the culturally meaningful parts of the reference or that a designer would prefer the constrained image.

The most immediate alternative explanation for the high structural score is objective alignment. The method is optimized to match reference edges, and the reported score measures closely related edge agreement. This is expected behavior rather than an invalid computation, but it places the result at the level of an optimization diagnostic. A method could raise the score by imposing reference layout at the expense of carrier-specific placement. Independent motif annotation, recognition judgments, and carrier-layout assessment would distinguish useful structural retention from this competing explanation.

The observed content-loss increase further qualifies the phrase ``without sacrificing content.'' Its aggregate magnitude is small relative to the reported content-loss scale, yet there is no prespecified tolerance, equivalence test, or perceptual calibration that would make the change negligible by definition. The correct interpretation is a measured trade-off between the edge target and the content objective. Reporting that trade-off allows later studies to select a design-appropriate operating point rather than treating one coefficient as universally optimal.

\subsection{Control Mechanisms and Prior Work}
Separately weighted objectives make a synthesis procedure inspectable: a designer or researcher can identify which term was changed and examine the resulting image. This operational property is distinct from semantic disentanglement. Earlier perceptual-factor control and pyramid-based stylization already address spatial, color, and scale-related appearance~\cite{gatys2017control,lapstyle2021}. The present multi-scale variant should be understood in that context, as a fixed-filter combination studied on product images, not as the first decomposition of visual style.

The single intervention case cannot determine whether a term's influence is restricted to its label. The color-named term depends on full-image features, while middle-band edges and high-band variance respond to multiple visual attributes. A useful next experiment would therefore define observable attributes before changing weights, evaluate intended and unintended changes on every pair, and include simpler controls. This would test whether the additional terms provide usable selectivity beyond merely producing visibly different images.

The local covariance and attention results also require implementation-aware interpretation. Matching their targets by iterative reconstruction differs from using the original trained decoders and objectives. The results show that the reference-edge term can act in those local pipelines; they do not resolve whether the same trade-off holds in official implementations with matched preprocessing and runtime. That question requires a separate reproduction study rather than a broader claim inferred from the alignment mechanism alone.

\subsection{Cultural Interpretation and Independent Measurement}
A cultural reference has more than a contour. Composition, symbolism, context, provenance, and the appropriateness of reuse can affect the interpretation of a product visualization. Neither Sobel edges nor a relative Gram-distance ratio represents these factors explicitly. Culturally situated assessment, such as the dialogue-based approach discussed by Qadri et al.~\cite{qadri2024}, makes clear why technical similarity should be combined with accountable domain interpretation rather than relabeled as cultural authenticity.

The distinction changes how CFM should be used. The objective components are useful for diagnosing a synthesis and identifying examples for systematic review. They are not a validated cultural-fidelity scale, and they should not be used to select only favorable images for a human study. An independent evaluation would present the reference and carrier, randomize method identities, collect judgments under a documented consent and recruitment procedure, and preserve all planned comparisons. Symbol recognizability, appropriateness, and overall aesthetics should be treated as distinct endpoints. Associations with the training-related score would then remain exploratory until replication supports broader validity.

The current study does not supply those human observations. Prepared instruments and simulated ratings can support software checks, but cannot establish cultural or perceptual conclusions. Removing numerical human-validation claims is therefore necessary to keep the manuscript aligned with its available evidence; it is not a substitute for collecting the missing evidence.

\subsection{Reproducibility and Generalization Limits}
Three implementation issues materially bound the computational findings. First, the archived preprocessing applies ImageNet normalization twice. Although the paired base comparisons share this pipeline, the deviation changes feature scales and complicates comparisons with standard VGG-based methods. Correcting it would require new, explicitly versioned paired runs rather than reinterpretation of the existing numbers. Second, the saved VGG-distance table covers only one carrier because of filename parsing. Fixing that evaluation is necessary for full benchmark coverage, and its present subset cannot establish a universal no-cost trade-off. Third, rounded style-loss values and missing per-run termination and runtime records limit precision and efficiency claims.

The sampling design introduces a separate limit. Five references and three ceramic tea-ware images provide a small, crossed pilot collection with shared sources across pairs. They do not characterize the diversity of each cultural tradition or transfer to textiles, jewelry, architecture, or other product materials. The multi-scale evidence is narrower still, covering one pair. Expanding the study should separate development from evaluation by source image and cultural-reference family, retain unsuccessful outcomes, and assess sensitivity to masks, crop alignment, and resolution.

Data rights and provenance also require documentation before public release. The inspected code and results identify operational filenames, but the present record does not establish a redistributable license and cultural-source history for every reference and carrier. Releasing metadata and code can improve traceability, but does not itself confer image permissions. These constraints affect reproducibility and responsible reuse, even when the computations are technically repeatable.

\section{Conclusion}
\label{sec:conclusion}
We examined explicit reference-edge constraints and separately weighted multi-scale objectives for stylizing ceramic product images. On the available 15-pair collection, the edge constraint increased reference-edge agreement while slightly increasing the mean content objective. Additional local alignment controls showed compatible diagnostic changes, and a single loss-removal case illustrated the multi-scale interface. The results support a transparent mechanism for controlling a reference-derived visual prior; they do not establish independent cultural-semantic validity, universally independent visual layers, or a general advantage over published stylization systems. Standard-preprocessing replications, complete diagnostic coverage, and independent cultural and perceptual evaluation are required to advance the framework beyond this pilot scope.

\section*{Data and Code Availability}
The revision archive contains computational scripts, saved tabular results, summary scripts, and figure provenance. A public repository and a reproducible software environment have not been established. Image redistribution and checkpoint use remain subject to source and permission verification. Simulated human-rating records are excluded from empirical analyses.
